\documentclass[10pt]{article}
\usepackage[margin=1in]{geometry}
\usepackage{amsmath,amssymb,booktabs,hyperref}
\usepackage[numbers,sort&compress]{natbib}
\title{Big Fast Weights, Small Slow Weights: Meta-Learning a LoRA Initialization for End-to-End Test-Time Training \\ \large Preprint proposal (draft, planning stage)}
\author{Helen Hui}
\date{September 2026}
\begin{document}
\maketitle

\section{Motivation}
TTT-E2E \citep{tandon2025e2e} treats long-context language modeling as continual learning: a Transformer with sliding-window attention (window $k{=}8$K) keeps learning at test time by taking one gradient step of the next-token loss per 1K-token chunk, and its initialization is meta-learned so that the loss \emph{after} test-time training is what pre-training minimizes. In that design the fast (inner-loop) weights are deliberately small: a second MLP in the last quarter of the blocks, about 10\% of parameters, while every parameter, including the fast-weight initialization, is a slow (outer-loop) weight. Two observations motivate inverting this allocation. First, the paper's own ablation shows that context scaling improves with the amount of fast state (updating 6 or 12 of 24 layers scales like full attention; 1 or 3 does not), and LaCT reports gains from fast weights that reach 40\% of the model \citep{zhang2025lact}. Second, TTT-E2E's outer loop is expensive (gradients of gradients, $3.4\times$ slower than standard pre-training at 8K), and the authors list ``initialize from a pre-trained Transformer'' as the way to amortize it. If a pre-trained model's own MLP matrices can serve as large fast weights, then the outer loop only needs to learn how to \emph{use} them, and a small parameter set may suffice for that. We therefore ask: on a pre-trained sliding-window Transformer, can a small meta-learned slow set, a LoRA on the attention projections \citep{hu2022lora} plus normalization gains and per-tensor inner learning rates, recover most of the benefit of full meta-learning when the fast weights are the model's full MLPs?

A caution shapes the design. A LoRA placed on the same MLP that the inner loop updates is only a low-rank shift of the fast-weight initialization; it cannot change what the attention layers write into fast memory, which is what an all-weights outer loop learns. The slow set is therefore centered on attention, with MLP-LoRA kept as an ablation. Overwriting a pre-trained MLP also removes the static ``safe storage'' MLP that TTT-E2E keeps to protect pre-trained knowledge, so forgetting is measured directly.

\section{Related work}
\textbf{Test-time training for sequences.} TTT layers \citep{sun2024ttt} update a per-layer fast weight with a self-supervised reconstruction loss and use mini-batch inner steps with gradient checkpointing through time; TTT-E2E \citep{tandon2025e2e} replaces the layer-wise loss with the end-to-end next-token loss and meta-learns all weights. LaCT \citep{zhang2025lact} uses large SwiGLU fast weights with 2K to 1M-token chunks, weight normalization, and optionally Muon-style orthogonalized updates \citep{liu2025muon}. Titans \citep{behrouz2024titans} adds momentum and a data-dependent forgetting gate to the fast-weight update. In-Place TTT \citep{feng2026inplace} updates the MLP down-projection of a pre-trained model with first-order steps and no outer loop. Dynamic evaluation \citep{krause2018dynamic} is test-time gradient descent on all weights without meta-learning and is our TTT-naive baseline. Clark et al.\ \citep{clark2022fastweight} meta-learn a fast-weight MLP appended to a full-attention Transformer.

\textbf{Meta-learning and low-rank adapters.} MAML \citep{finn2017maml} learns an initialization through gradients of gradients; Meta-SGD \citep{li2017metasgd} and MAML++ \citep{antoniou2019mamlpp} additionally learn per-parameter or per-layer inner learning rates, which we adopt as slow parameters. PERK \citep{chen2025perk} meta-learns a LoRA initialization whose inner loop also lives in the LoRA; FocuSFT \citep{pei2026focusft} is the inverse of our setting (fast LoRA, slow full weights, first-order). rsLoRA \citep{kalajdzievski2023rslora} gives the $\alpha/\sqrt{r}$ scaling needed for large ranks, and Biderman et al.\ \citep{biderman2024lora} show that continued pre-training requires higher ranks than instruction tuning. Memory-efficient meta-gradients via mixed-mode differentiation \citep{kemaev2025mixflow} target the second-order activation term rather than the per-step fast-weight copies that dominate our setting; pathologies of long unrolls are analyzed in \citep{metz2019pathologies}. To our knowledge no prior work meta-learns a LoRA-only slow set through an inner loop whose fast weights are the full MLPs of a pre-trained model.

\section{Method}
\textbf{Model.} Llama-3.2-1B (16 blocks, $d{=}2048$, MLP width 8192, GQA with 8 KV heads) loaded into the TTT-E2E JAX codebase, with every attention layer run as sliding-window attention with $k{=}8192$; at the 8K pre-training context this is identical to full attention. Fast weights $W$ are the MLP matrices $(W_1,W_2,W_3)$ of the last $L_f\in\{4,8,16\}$ blocks (201M, 403M, 805M parameters). Slow weights $\theta$ are: LoRA adapters $B_iA_i$ on $W_q,W_k,W_v,W_o$ of all blocks (rank $r\in\{16,64,256\}$, rsLoRA scale $\alpha/\sqrt r$, $\alpha{=}16$, $B{=}0$ at initialization), all RMSNorm gains, and one learned log-learning-rate per fast tensor. Everything else is frozen at its pre-trained value. Fast weights are reset to the pre-trained values at every sequence boundary.

\textbf{Inner loop.} A sequence of $T$ tokens is split into $N{=}T/b$ chunks with $b{=}1024$. For chunk $i$ the loss $\ell_i(W_{i-1})$ is the next-token cross-entropy on the chunk with the current fast weights, and one step gives $W_i$. Two update rules are compared. Normalized SGD, applied per tensor,
\[
W_i = W_{i-1} - \eta\, e^{\vartheta}\sqrt{mn}\;\frac{g_i}{\|g_i\|_F+\varepsilon},\qquad g_i=\nabla_W \ell_i(W_{i-1}),
\]
where $\eta$ is the per-element RMS step size and $\vartheta$ the learned per-tensor scalar; and AdamW with $\varepsilon{=}10^{-8}$, moments warm-started from the first chunk's gradient and carried in fp32. Optionally $W_{i-1}$ is decayed toward $W_0$ by a factor $(1-\lambda)$ before each step. The outer loss is the paper's $\mathcal L(\theta)=\frac1T\sum_i \sum_{t\in\text{chunk }i}\ell_t(W_{i-1})$, differentiated through the whole inner trajectory.

\textbf{Outer loop.} AdamW ($\beta{=}(0.9,0.95)$, weight decay $0.1$ on LoRA only, gradient clip $1.0$, 10\% warm-up, cosine decay), 0.5M tokens per step, learning rate swept over $\{3\times10^{-5},\dots,10^{-2}\}$. Stage 1 meta-trains at 8K on DCLM (documents $\ge$8K tokens) for 1.3B tokens, 5\% of the model's Chinchilla budget, mirroring the paper's fine-tuning recipe; Stage 2 extends to 32K on Books with the same budget.

\textbf{Memory.} With checkpointing through time in groups of $n$ chunks, the fast-weight copies held per sequence are $(N/n+n)$, each of size $|W|$ plus optimizer state; with $n\approx\sqrt N$ this is 9 GB (normalized SGD) or 27 GB (AdamW) at 32K for $L_f{=}4$, and 36 GB or 109 GB for $L_f{=}16$, the latter requiring fast-weight sharding or host offload.

\textbf{Arms and metrics.} (A) frozen SWA model; (B) TTT-naive, inner loop only; (C) proposed; (D) all weights slow; (E) TTT-E2E as released; (F) the paper's small prime-MLP fast weights under the same pre-trained initialization and budget as (C), which isolates fast-weight size from the training regime. Metrics: held-out loss on DCLM (8K) and Books (32K); loss by token index; a forgetting probe (change in loss on an unrelated held-out chunk after test-time training, relative to $W_0$); peak device memory and training seconds per 1K tokens; deployed parameter count after merging the LoRA. The hypothesis is supported if (C) approaches (D) and both exceed (B); it is falsified if (C) matches (B) while (D) does not.

\bibliographystyle{plainnat}
\bibliography{refs}
\end{document}
